\documentclass[11pt,a4paper]{article}

\usepackage[margin=1in]{geometry}
\usepackage{xcolor}
\usepackage{parskip}
\usepackage{enumitem}
\usepackage{hyperref}
\usepackage{amsmath,amssymb}
\usepackage{bm}
\usepackage{siunitx}
\usepackage[authoryear,round]{natbib}
\usepackage{xr-hyper}

\hypersetup{
  colorlinks=true,
  linkcolor=blue!60!black,
  urlcolor=blue!60!black,
  citecolor=blue!60!black,
}

% Pull labels (sections, figures, tables, equations) from the main
% manuscript's .aux file so that \ref{...} resolves to the numbering
% in the compiled manuscript. Requires main.tex to have been built first.
\makeatletter
\let\coverletter@orig@bibcite\bibcite
\let\bibcite\@gobbletwo
\externaldocument{../../build/main}
\let\bibcite\coverletter@orig@bibcite
\makeatother

\newcommand{\Kzero}{\ensuremath{K_0}}

\title{Cover Letter\\
  \large Manuscript COMGE-D-26-01109 (R1)}
\author{Yusong Han, Ryosuke Uzuoka, Kyohei Ueda}
\date{\today}

\begin{document}

\maketitle

Dear Editor,

Thank you for the opportunity to revise our manuscript
``\emph{Influence of Initial Stress Anisotropy on Liquefaction Resistance
in Undrained Cyclic Torsional Shear: DEM Simulation of Hollow Cylinder
Tests}'' (COMGE-D-26-01109) submitted to
\emph{Computers and Geotechnics}. We sincerely thank both reviewers for their
thoughtful and constructive comments, which have helped us to substantially
improve the manuscript.

We have carefully addressed every comment. A detailed point-by-point response
is provided in the accompanying response-to-reviewers document. In the revised
manuscript, added or modified text is marked using \texttt{latexdiff} against
the submitted version (also provided).

\section*{Summary of major revisions}

\begin{itemize}[leftmargin=*]
  \item \textbf{Clarified the relative-density labels}
    ``$D_r \approx 90\%$'' and ``$D_r \approx 75\%$'' as
    behaviour-matched references to the Toyoura density states
    reported by \citet{Nakata1998} and \citet{Ishihara1985}, rather
    than values computed from the DEM void ratio via the textbook
    formula.
  \item \textbf{Expanded the Introduction} with a review of prior
    DEM-HCA models and the boundary-condition taxonomy (rigid wall,
    stacked wall, flexible membrane); the rigid-wall configuration of
    the present work is positioned as the required DEM realisation of
    the undrained constant-volume boundary.
  \item \textbf{Added a consolidated case table}
    (Table~\ref{tab:case_matrix}) enumerating the CSR levels tested
    at each $(D_r, \Kzero)$ combination across the primary, expanded,
    and verification sets.
  \item \textbf{Added a new ``Numerical verification'' subsection}
    (Section~\ref{subsec:numerical_verification},
    Fig.~\ref{fig:numerical_verification}) reporting the volume
    drift, the radial and axial servo-tracking errors, and the
    inertial number for a representative post-liquefaction run,
    documenting both exact volume conservation and the quasi-static
    character of the loading.
  \item \textbf{Identified the source of the CSL in
    Fig.~\ref{fig:macro_response_c}} as the steady-state friction
    angle $\varphi_{cs}\approx 31^{\circ}$ reported by
    \citet{VerdugoIshihara1996} for Toyoura sand, converted to the
    HCA pure-shear loading ($b=0.5$) via the Mohr--Coulomb relation
    $M_{cs}=\sqrt{3}\sin\varphi_{cs}\approx 0.892$; the citation and
    conversion are now stated in the figure caption.
  \item \textbf{Added contact-force anisotropy analysis} to complement
    the existing contact-density analysis: a new normal-force rose
    (Fig.~\ref{fig:contact_force_sequence}) and a dual-axis upgrade
    to Fig.~\ref{fig:fabric_mechanism} that overlays the count- and
    force-weighted torsion-resisting fractions across the full
    expanded dataset (10 cases), supported by a new appendix table
    (Table~\ref{tab:fabric_diagonal}).
  \item \textbf{Quantified the differences} between specimens in
    Fig.~\ref{fig:force_disp_sequence} via the new normal-force rose
    (Fig.~\ref{fig:contact_force_sequence}), which reports the
    force-fabric deviator $A_n$, principal-direction tilt $B_n$, and
    specimen-mean $\langle F_n\rangle$ for each panel.
  \item Minor clarifications and additional references throughout.
\end{itemize}

We hope the revised manuscript now meets the journal's standards and the
reviewers' expectations. We remain grateful for the reviewers' careful
reading and valuable suggestions.

\vspace{1em}

\noindent Sincerely,

\vspace{0.5em}

\noindent Yusong Han, Ryosuke Uzuoka, Kyohei Ueda \\
on behalf of all authors

\bibliographystyle{plainnat}
\bibliography{refs}

\end{document}
